\subsection{THE AXIOM OF THE ORDERED PAIR}

As we have said in §1, no. 1, the sign \(\supset\) is a substantific sign of weight 2 in the theory of sets. If T, U are terms, then \(\supset\) TU is a term, which in practice is denoted by (T, U). In this notation, the axiom of the ordered pair is the following :

\[
(\forall x) (\forall x ^ {\prime}) (\forall y) (\forall y ^ {\prime}) (((x, y) = (x ^ {\prime}, y ^ {\prime})) \Longrightarrow (x = x ^ {\prime} \quad \text { and } \quad y = y ^ {\prime})). \tag {A3.}
\]

Since the relation `` \(x = x'\) and \(y = y'\)'' implies \((x, y) = (x', y')\) by C44 (Chapter I, § 5, no. 2), the relation \((x, y) = (x', y')\) is equivalent to `` \(x = x'\) and \(y = y'\)''.

¶ The relation \((\exists x)(\exists y)(z=(x,y))\) is written ``z is an ordered pair''. If z is an ordered pair, the relations

\[
(\exists y) (z = (x, y)), \quad (\exists x) (z = (x, y))
\]

are functional with respect to \(x\) and \(y\) respectively; this is an immediate consequence of A3.

The terms

\[
\tau_ {x} ((\exists y) (z = (x, y))) \quad \text { and } \quad \tau_ {y} ((\exists x) (z = (x, y)))
\]

are denoted by \(\mathbf{pr}_1z\) and \(\mathbf{pr}_2z\) , respectively, and are called the first coordinate (or first projection) and the second coordinate (or second projection) of \(z\) . If \(z\) is an ordered pair, the relation \((\exists y)(z = (x, y))\) is equivalent to \(x = \mathbf{pr}_1z\) , and the relation \((\exists x)(z = (x, y))\) to \(y = \mathbf{pr}_2z\) (Chapter I, §5, no. 3).

The relation \(z = (x, y)\) is equivalent to `` \(z\) is an ordered pair, and \(x = \text{pr}_1z\) and \(y = \text{pr}_2z\) ''; for the latter relation is equivalent to

\[
(\exists x ^ {\prime}) (\exists y ^ {\prime}) (\exists x ^ {\prime \prime}) (\exists y ^ {\prime \prime}) (z = (x ^ {\prime}, y ^ {\prime}) \text {   and   } z = (x, y ^ {\prime \prime}) \text {   and   } z = (x ^ {\prime \prime}, y));
\]

by A3, `` \(z = (x', y')\) and \(z = (x, y'')\) and \(z = (x'', y)\)'' is equivalent to `` \(z = (x, y)\) and \(x = x'\) and \(x = x''\) and \(y = y'\) and \(y = y''\) ''; hence by C33 (Chapter I, §4, no. 3), `` \(z\) is an ordered pair, and \(x = \text{pr}_1 z\) and \(y = \text{pr}_2 z\) '' is equivalent to

\[
z = (x, y) \text { and } (\exists x ^ {\prime}) (\exists y ^ {\prime}) (\exists x ^ {\prime \prime}) (\exists y ^ {\prime \prime}) (x = x ^ {\prime} \text { and } x = x ^ {\prime \prime} \text { and } y = y ^ {\prime} \text { and } y = y ^ {\prime \prime}),
\]

which proves our assertion. Evidently we have

\[
\operatorname{pr} _ {1} (x, y) = x, \quad \operatorname{pr} _ {2} (x, y) = y,
\]

and the relation \(z = (\mathrm{pr}_1(z), \mathrm{pr}_2(z))\) is equivalent to `` \(z\) is an ordered pair''.

Let \(R\{x, y\}\) be a relation, where the letters x and y are distinct and appear in R. Let z be a letter, distinct from x and y, which does not appear in R. Let \(S\{z\}\) denote the relation

\[
(\exists x) (\exists y) (z = (x, y) \text {   and   } R \{x, y \});
\]

this is a relation which contains one letter fewer than R, and which is equivalent to ``z is an ordered pair and \(R\{pr_{1}z, pr_{2}z\}\) '', this follows from the fact that \(z = (x, y)\) is equivalent to ``z is an ordered pair and \(x = pr_{1}x\) and \(y = pr_{2}z\) '', and from criteria C33 (Chapter I, §4, no. 3)

and C47 (Chapter I, §5, no. 3). It follows immediately that \(R\{x, y\}\) is equivalent to \(S\{(x, y)\}\) , and also to

\[
(\exists z) (z = (x, y) \text {   and   } S \{z \})
\]

by C47.

This means that we may interpret a relation between the objects x and y as a property of the ordered pair formed by these objects.

